\section{Un contrat d'assurance}
\label{sec:U6}

Les cinq unités précédentes ont introduit, une par une, toutes les pièces
nécessaires : l'actualisation, les dates multiples, la dépendance aux taux de
marché, l'aléa, l'optionalité. Il n'y a rien de neuf à apprendre ici. On les
assemble.

Le résultat porte un nom réglementaire --- la meilleure estimation --- mais il
n'est rien d'autre que la valeur actuelle probable d'un échéancier. Cette unité
est l'atterrissage de toute la progression.

\subsection{La notion : l'échéancier probabilisé d'un contrat}

\begin{notion}{Un contrat est une variable aléatoire}
Un contrat d'assurance engage l'assureur à verser des prestations dont le
montant et la date sont incertains. Il est donc décrit par une suite de flux
\emph{attendus} $(\mathbb{E}[F_1],\dots,\mathbb{E}[F_n])$, et sa valeur est
\[
  \mathrm{BE} \;=\; \sum_{k=1}^{n} \frac{\mathbb{E}[F_k]}{(1+z_k)^{k}} .
\]
Chaque ingrédient a été vu séparément : l'espérance vient de l'unité~4, les
dates multiples de l'unité~2, les taux $z_k$ lus sur la courbe de l'unité~3.

\medskip
En assurance non-vie, les flux attendus se déduisent d'une \emph{cadence de
règlement}\index{cadence de règlement} : la charge ultime $S$ d'un exercice de
survenance se règle progressivement selon un profil $c_1,\dots,c_n$ observé sur
l'historique, avec $\sum_k c_k = 1$. Alors $\mathbb{E}[F_k] = S \times c_k$.
\end{notion}

\begin{center}
\begin{tikzpicture}[x=1.15cm, y=0.75cm]
  \draw[timeaxis] (-0.4,0) -- (5.9,0);
  \foreach \x/\l in {0/0, 1/1, 2/2, 3/3, 4/4, 5/5}
    { \node[timenode] at (\x,0) {}; \node[etiq, below=3pt] at (\x,0) {\l}; }
  \draw[flux] (1,0) -- (1,1.75); \node[etiq, above] at (1,1.75) {$Sc_1$};
  \draw[flux] (2,0) -- (2,1.30); \node[etiq, above] at (2,1.30) {$Sc_2$};
  \draw[flux] (3,0) -- (3,0.95); \node[etiq, above] at (3,0.95) {$Sc_3$};
  \draw[flux] (4,0) -- (4,0.65); \node[etiq, above] at (4,0.65) {$Sc_4$};
  \draw[flux] (5,0) -- (5,0.40); \node[etiq, above] at (5,0.40) {$Sc_5$};
  \draw[fluxneg] (0,1.55) -- (0,0.05);
  \node[etiq, above] at (0,1.6) {$\mathrm{BE}$};
\end{tikzpicture}
\end{center}

\begin{exercice}[meilleure estimation en non-vie]
Un portefeuille de responsabilité civile automobile présente une charge ultime
estimée à $8\,000\,000$~\euro{} au titre des sinistres déjà survenus. La cadence
de règlement observée est de $40\,\%$, $25\,\%$, $18\,\%$, $12\,\%$ et $5\,\%$,
les règlements intervenant en fin d'année. Le taux d'actualisation est de
$2{,}20\,\%$, constant.

Calculer la meilleure estimation de la provision pour sinistres à payer, puis la
comparer à la provision non actualisée.
\end{exercice}

\begin{corrige}
\begin{center}\small
\begin{tabular}{crrrr}
\toprule
Année & Part réglée & Flux attendu & Facteur & Flux actualisé \\
\midrule
1 & $40\,\%$ & $3\,200\,000$ & $0{,}97847$ & $3\,131\,115$ \\
2 & $25\,\%$ & $2\,000\,000$ & $0{,}95741$ & $1\,914\,821$ \\
3 & $18\,\%$ & $1\,440\,000$ & $0{,}93680$ & $1\,348\,993$ \\
4 & $12\,\%$ & $960\,000$    & $0{,}91663$ & $879\,970$ \\
5 & $5\,\%$  & $400\,000$    & $0{,}89690$ & $358\,761$ \\
\midrule
\multicolumn{3}{r}{\textbf{Meilleure estimation}} & & $\mathbf{7\,633\,661}$ \\
\bottomrule
\end{tabular}
\end{center}
La provision non actualisée s'élève à $8\,000\,000$~\euro. L'écart, soit
$366\,339$~\euro{} ou $4{,}58\,\%$, est le seul effet de l'actualisation.

\medskip
Cet écart dépend de deux choses : le niveau des taux, et la longueur de la
cadence. Sur une branche à déroulement long --- la responsabilité civile
médicale, la construction --- où les règlements s'étalent sur quinze ou vingt
ans, il peut dépasser $15\,\%$.
\end{corrige}

\begin{exercice}[l'écart avec la provision comptable]
Le même portefeuille est provisionné en normes françaises à $8\,600\,000$~\euro,
la réglementation imposant une évaluation prudente et sans actualisation.

Décomposer l'écart entre cette provision et la meilleure estimation calculée
ci-dessus, en distinguant ses deux sources.
\end{exercice}

\begin{corrige}
\[
  8\,600\,000 - 7\,633\,661 = 966\,339\ \text{\euro}.
\]
Cet écart se décompose exactement en deux termes :
\begin{center}\small
\begin{tabular}{lrl}
\toprule
Source & Montant & Origine \\
\midrule
Marge de prudence   & $600\,000$ & hypothèses prudentes contre réalistes \\
Actualisation       & $366\,339$ & valeur temporelle de l'argent \\
\midrule
\textbf{Écart total} & $\mathbf{966\,339}$ & \\
\bottomrule
\end{tabular}
\end{center}

\medskip
Les deux sources sont de nature différente et se justifient devant un contrôleur
par des arguments distincts. La marge de prudence relève du choix des hypothèses,
étudié à l'unité~4 ; l'actualisation relève de la mécanique financière, étudiée
à l'unité~2.

Le passage à la meilleure estimation \emph{libère} donc des fonds propres --- ici
près d'un million d'euros. C'est l'un des effets les plus visibles du passage
en normes prudentielles, et l'une des questions les plus fréquentes en contrôle.
\end{corrige}

\subsection{Généralisation : frontière des contrats et marge de risque}

\begin{generalisation}{Ce que l'on projette, et ce que l'on ajoute}
\textbf{La frontière des contrats.} Tous les flux futurs ne sont pas à
provisionner. La \emph{frontière du contrat}\index{frontière des contrats}
délimite les engagements que l'assureur ne peut plus refuser : les primes
futures ne sont incluses que s'il ne dispose ni du droit de résilier, ni de
celui de refuser la prime, ni de celui d'en modifier librement le niveau. Un
contrat annuel tacitement reconductible, résiliable à l'échéance, a donc une
frontière d'un an --- même si le client restera vraisemblablement dix ans.

C'est une frontière \emph{juridique}, pas comportementale. Elle produit des
résultats contre-intuitifs, et elle est la source d'erreur la plus fréquente
dans un calcul de meilleure estimation vie.

\medskip
\textbf{La meilleure estimation est nette.} On projette tous les flux liés au
contrat : prestations, frais de gestion et d'acquisition, et primes à recevoir
en signe opposé.
\[
  \mathrm{BE} \;=\; \sum_{k} \frac{\mathbb{E}[\text{prestations}_k] + \mathbb{E}[\text{frais}_k]
    - \mathbb{E}[\text{primes}_k]}{(1+z_k)^{k}} .
\]
Une meilleure estimation peut donc être négative, ce qu'aucune provision
comptable ne permet.

\medskip
\textbf{La marge de risque.} La meilleure estimation est une espérance : par
construction, elle sera insuffisante une fois sur deux. Les provisions
techniques y ajoutent une \emph{marge de risque}\index{marge de risque},
définie comme le coût d'immobilisation du capital qu'un repreneur exigerait
pour porter les engagements jusqu'à leur extinction :
\[
  \mathrm{MR} \;=\; \mathrm{CoC} \times \sum_{k \geq 0}
    \frac{\mathrm{SCR}_k}{(1+z_{k+1})^{k+1}},
  \qquad \mathrm{CoC} = 6\,\% .
\]
On retrouve, une fois encore, une somme de flux actualisés --- les $\mathrm{SCR}_k$
étant les capitaux requis projetés année après année.

\medskip
\[
  \boxed{\ \text{Provisions techniques} \;=\; \mathrm{BE} \;+\; \mathrm{MR}\ }
\]
\end{generalisation}

\begin{exercice}[frontière des contrats]
Pour chacune des situations suivantes, indiquer si les primes futures entrent
dans la meilleure estimation, et justifier.

\textbf{(a)} Un contrat automobile annuel, tacitement reconductible, que
l'assureur peut résilier à chaque échéance.

\textbf{(b)} Un contrat d'épargne à versements libres, sans engagement de
l'assuré ni faculté de refus par l'assureur.

\textbf{(c)} Un contrat de prévoyance décès à prime nivelée sur vingt ans, que
l'assureur ne peut ni résilier ni retarifer.
\end{exercice}

\begin{corrige}
\textbf{(a)} \emph{Non.} La faculté de résiliation annuelle place la frontière à
l'échéance suivante. Le fait que le contrat sera vraisemblablement reconduit dix
ans est sans effet : le critère est juridique, non comportemental.

\medskip
\textbf{(b)} \emph{Non} pour les versements futurs, qui ne sont pas des
engagements : l'assuré est libre de ne pas verser. En revanche, l'épargne déjà
constituée est intégralement provisionnée, avec toutes les garanties qui s'y
attachent.

\medskip
\textbf{(c)} \emph{Oui}, sur vingt ans. L'assureur est engagé pour toute la
durée sans pouvoir modifier le tarif : les primes futures sont dans la frontière
et viennent en déduction de la meilleure estimation.

\medskip
Le cas (c) illustre la possibilité d'une meilleure estimation négative : si la
valeur actuelle des primes excède celle des prestations et des frais, le contrat
constitue un actif, non un passif. C'est économiquement juste et comptablement
impossible en normes françaises.
\end{corrige}

\begin{exercice}[provisions techniques complètes]
Reprendre le portefeuille de l'exercice~6.1, dont la meilleure estimation
s'établit à $7\,633\,661$~\euro. Le capital de solvabilité requis au titre de ce
portefeuille est projeté comme suit, en fin de chaque année :
\begin{center}\small
\begin{tabular}{lrrrrr}
\toprule
Année $k$ & $0$ & $1$ & $2$ & $3$ & $4$ \\
$\mathrm{SCR}_k$ & $1\,200\,000$ & $780\,000$ & $460\,000$ & $220\,000$ & $80\,000$ \\
\bottomrule
\end{tabular}
\end{center}
Le taux d'actualisation reste de $2{,}20\,\%$ et le coût du capital de $6\,\%$.

Calculer la marge de risque, puis les provisions techniques.
\end{exercice}

\begin{corrige}
Chaque $\mathrm{SCR}_k$ est actualisé sur $k+1$ périodes :
\begin{center}\small
\begin{tabular}{crrr}
\toprule
$k$ & $\mathrm{SCR}_k$ & Facteur $(1{,}022)^{-(k+1)}$ & Montant actualisé \\
\midrule
0 & $1\,200\,000$ & $0{,}97847$ & $1\,174\,168$ \\
1 & $780\,000$    & $0{,}95741$ & $746\,780$ \\
2 & $460\,000$    & $0{,}93680$ & $430\,928$ \\
3 & $220\,000$    & $0{,}91663$ & $201\,660$ \\
4 & $80\,000$     & $0{,}89690$ & $71\,752$ \\
\midrule
\multicolumn{3}{r}{\textbf{Somme}} & $\mathbf{2\,625\,289}$ \\
\bottomrule
\end{tabular}
\end{center}
\[
  \mathrm{MR} = 6\,\% \times 2\,625\,289 = 157\,517\ \text{\euro}.
\]
\[
  \text{Provisions techniques} = 7\,633\,661 + 157\,517 = 7\,791\,178\ \text{\euro}.
\]

\medskip
La marge de risque représente ici $2{,}06\,\%$ de la meilleure estimation. Sur
une branche à déroulement long, où les $\mathrm{SCR}$ projetés restent élevés
pendant quinze ans, cette proportion est sensiblement plus forte.

Noter que même après ajout de la marge de risque, les provisions techniques
restent inférieures de $808\,822$~\euro{} à la provision comptable de
$8\,600\,000$~\euro.
\end{corrige}

\begin{pont}{Où cela sert dans Solvabilité~II}
Cette unité \emph{est} le passif du bilan prudentiel. Rien de plus n'est requis
pour le construire, du moins en assurance non-vie.

\medskip
La progression complète se relit d'un trait. Un flux certain, actualisé
(unité~1). Plusieurs dates, donc une somme (unité~2). Des taux lus sur la courbe
EIOPA plutôt qu'un taux plat (unité~3). Des flux pondérés par leur probabilité
d'occurrence (unité~4). Une pondération risque-neutre lorsque le contrat comporte
des options (unité~5). Enfin l'assemblage, augmenté de la marge de risque
(unité~6). Chaque étage n'ajoutait qu'un ingrédient, et le dernier ne fait que
les réunir.

\medskip
Deux prolongements restent à construire, et ils forment la suite du cours.

\emph{L'actif face au passif.} Une meilleure estimation ne se calcule pas
isolément : sa valeur dépend du rendement des placements qui la couvrent, et
réciproquement les décisions d'investissement dépendent du profil des
engagements. Cette interaction est l'objet de la gestion actif-passif, à
l'unité~7.

\emph{La projection stochastique.} En assurance vie, les options de l'unité~5
--- garanties de taux, participation aux bénéfices, rachats --- interdisent la
projection déterministe employée ici. Il faut simuler des milliers de
trajectoires et moyenner : c'est l'unité~8.

\medskip
La méthode reste pourtant la même dans les deux cas : une espérance de flux
actualisés. Seule la manière de calculer l'espérance change.
\end{pont}

\begin{refreg}
\textbf{Références.} Directive 2009/138/CE, article~76 : les provisions
techniques correspondent au montant que les entreprises devraient payer pour
transférer immédiatement leurs engagements à une autre entreprise. \\
Article~77 : la valeur des provisions techniques est égale à la somme de la
meilleure estimation et de la marge de risque, cette dernière étant calculée de
manière à garantir que la valeur des provisions équivaut au montant exigé pour
reprendre et honorer les engagements. \\
Le coût du capital retenu pour la marge de risque est fixé réglementairement. \\
Règlement délégué (UE) 2015/35 : frontière des contrats, segmentation et
méthodes de calcul de la marge de risque. \emph{[Références d'articles et taux
de coût du capital à vérifier avant diffusion.]}
\end{refreg}

\subsection*{Notations introduites}
\begin{description}[leftmargin=2.6cm, style=nextline, font=\normalfont]
  \notation{S}{charge ultime d'un exercice de survenance}
  \notation{c_k}{part de la charge réglée en année $k$}
  \notation{\mathbb{E}[F_k]}{flux attendu de la période $k$}
  \notation{\mathrm{BE}}{meilleure estimation}
  \notation{\mathrm{MR}}{marge de risque}
  \notation{\mathrm{CoC}}{coût du capital}
\end{description}
